The United States has about 4.2 million miles of public road, and roughly three-quarters of it---some 3.2 million miles---is owned not by states but by local governments \citep{fhwa2022}. Keeping this network drivable is expensive and chronically underfunded: state and local governments spend roughly \$206 billion a year on highways and roads, yet only 44\% of it goes to operations and maintenance rather than new capacity \citep{urban2021highway}. The bill for deferring upkeep is mounting. The 2025 ASCE Infrastructure Report Card grades the nation's existing roads a D+, with 39\% of roads in poor or mediocre condition and a projected \$684 billion funding gap over the next decade \citep{asce2025}, even as the federal gasoline tax that backstops the Highway Trust Fund has held at 18.4 cents per gallon since 1993 \citep{pgpf_htf_explained}. In a mature road network, the binding fiscal choice is increasingly not whether to build, but whether to maintain.

A great deal of existing research has focused on the effects of new roads\footnote{See \cite{ghani2016highway}, \cite{beenstock2016hedonic}, \cite{levkovich2016effects}, \cite{hoogendoorn2019house}, \cite{theisen2021road}, \cite{gibbons2019new}, \cite{fretz2022highways}. Some papers have also observed adverse effects of new roads such as state-led repression \citep{gonzalez2025} and environmental degradation \citep{asher2020eco}.}, especially in developing nations, providing valuable policy insights and spurring development initiatives like China's Belt and Road Initiative \citep{huang2016understanding} and India's Pradhan Mantri Gram Sadak Yojana \citep{asher2020}. However, the endogenous selection of these projects makes it difficult to separate the effects of new roads from the underlying economic conditions that motivated their construction, as infrastructure investments are not exogenous events but are deliberate policy choices. New road placement is selected by policymakers and may be endogenous to anticipated economic growth, unobserved local characteristics, or uncontrolled political factors. Moreover, in mature economies with extensive road networks, the central policy question is often not where to build new roads but how to maintain existing ones. 

% Local governments routinely ask residents to make fiscal choices whose costs are immediate but whose benefits accrue gradually, and road maintenance is a clear example. In Ohio, neighborhood roads are financed largely through dedicated property-tax levies that voters must renew every few years. A failed renewal does not block a new project; it removes an existing revenue stream that had been paying for ordinary resurfacing, patching, and preservation. The stakes are small for any one household but meaningful for a budget: in our sample the average road levy costs a household about \$76 per year, while the average failed renewal strips about \$163,547 from a local government---roughly 11\% of its road-maintenance budget. These are common, recurring events rather than rare megaprojects, which makes them both a routine public-management problem and a clean window onto a general question: what are the longer-run costs of choosing lower taxes over the continued upkeep of existing infrastructure?


% Local governments routinely ask residents to make fiscal choices about public services whose costs are immediate and whose benefits accrue gradually. Road maintenance is a clear example. A renewal levy can look like a small recurring tax bill for a household, but the revenue it preserves can determine whether a jurisdiction keeps up with ordinary resurfacing, patching, and preservation. When voters reject a renewal, the tax relief is visible right away. The infrastructure consequences are slower: maintenance is deferred, road quality can deteriorate, and the local housing market may eventually capitalize the decline in public-service quality.

% This paper studies that trade-off using local road-maintenance renewal elections in Ohio. Ohio local governments rely heavily on local property taxes to maintain neighborhood roads, and many dedicated road levies must be renewed by voters every five years. A failed renewal does not merely block a new project; it removes an existing revenue stream that had been financing ongoing maintenance. In our sample, the average road levy costs about \$76 per household per year and the average failed levy removes about \$163,547 from a local government, equal to roughly 11\% of its road-maintenance budget. These renewal votes therefore provide a useful setting for studying a common public-management question: what are the longer-run costs of choosing lower taxes over continued maintenance of existing infrastructure?

In this paper, we avoid endogeneity issues by studying property taxes that fund the maintenance of existing roads and by utilizing a quasi-experimental setting that arises from analyzing close elections.
We assemble a panel of 3,184 road-tax renewal referendums in Ohio from 1991 to 2021 and link the election records to financial records of local governments, CoreLogic house-sale transactions, community characteristics, and satellite imagery of local roads. Our empirical strategy compares jurisdictions where road-maintenance renewals narrowly fail with jurisdictions where they narrowly pass, using vote threshold to identify local fiscal shocks\footnote{\cite{grosz2025relaxing} study the effects of changes in these voting thresholds.} \citep{cellini2010value, biasi_what_2025}. The running variable is the share of votes against renewal, and treatment occurs when that share exceeds the 50\% threshold. Because renewal elections are scheduled by the expiration of prior levies rather than by a new discretionary investment proposal, close elections isolate local effects of losing an existing maintenance revenue source. Using our rich data, we first estimate the expected loss in the road maintenance budget for a local government after a tax cut from a failed renewal of road maintenance levy. Next, we estimate the effect of cutting local road maintenance taxes on housing prices. Finally, to understand the underlying channel, we fine-tune an Artificial Intelligence (AI) vision model using more than 53,000 satellite images from \cite{brewer2021} to predict changes in road quality after a road maintenance tax cut. 

Our empirical findings reveal that when a jurisdiction cuts its renewal road maintenance taxes, an average household saves about \$76 per year and the local government  faces an 11\% loss in maintenance funds, its house prices decrease by around \$15,350 (9\%) and road quality declines by up to 16\% in the precisely estimated post-election windows. In rolling 3-year post-election windows, we find little immediate road-quality response in years 1 to 3 after the vote, followed by statistically significant deterioration in years 2 to 4 and continued strong negative effects in years 3 to 5. A simple benefit-cost comparison shows that for each dollar of road-maintenance revenue voters forgo, households incur a net loss of \$1.4.

% After establishing the main results for house prices, we explore heterogeneity in the effects. We find that the effects on house prices are driven by urban areas rather than rural areas, with urban areas experiencing a 7.6\% decline in house prices while rural areas experience inconsistent changes in house prices. Moreover, we find an effect on the intensive margin with evidence for dosage-response based on the size of the levy. We also observe that higher-priced homes experience larger reductions in value relative to lower-priced homes. This suggests that wealthier households are more sensitive to road quality, as a poor road in front of a high-priced home is more likely to be noticed than a poor road in front of a low-priced home.

% The main result is that narrowly voting down a road-maintenance renewal is followed by a delayed decline in local house values. Median house prices are not statistically different around the cutoff before the vote, and effects remain small in the first few years afterward. Beginning four years after the vote, however, jurisdictions that narrowly fail to renew road levies have median house prices about \$15,350 lower on average over the following decade, or roughly 9\% of the average house value in the data. The timing is consistent with a maintenance mechanism: road conditions do not deteriorate immediately after a budget cut, but deferred upkeep becomes more visible over time.
  
We provide suggestive evidence on that mechanism using road imagery and a fine-tuned vision model that classifies road-surface quality. The imagery evidence shows no meaningful pre-election discontinuity in road quality. Post-election road-quality scores fall by 8 to 10 RQS points in the most precisely estimated delayed windows, roughly 13\% to 16\% relative to the pre-election mean. Taken together, the evidence is consistent with a channel in which failed renewals reduce available maintenance funding, which may lead to gradual declines in road quality and, over time, to reductions in house prices once deterioration becomes salient. 

The distribution of effects also matters for policy. The price declines are concentrated in urban jurisdictions and are larger for higher-value houses and larger levy failures. We interpret these patterns as evidence about incidence and capitalization rather than as direct evidence about household wealth or welfare. Sellers in affected markets may bear capitalized losses, buyers may benefit from lower purchase prices while inheriting lower service quality, and although current taxpayers receive short-run tax relief, these savings are eclipsed by the long-run decline in local asset values. The results therefore highlight this public finance trade-off and document the effects of public disinvestment decisions on local asset values.

The paper contributes to three strands of research. First, we extend work on infrastructure and housing markets from new construction toward the maintenance of existing public infrastructure. By exploiting close elections for renewal levies earmarked to fund road maintenance, we mitigate endogeneity concerns and estimate the capitalization of road maintenance decisions on local house prices. Second, we contribute to the local public finance literature by showing a striking example of salience-induced fiscal illusion, where voters in some local governments choose small but immediate tax savings at the cost of housing wealth destruction that only appears years later.
% how ballot-box decisions over dedicated revenue streams affect household wealth and highlighting the trade-off between short-term tax relief and long-term capital losses. 
Third, we demonstrate how scalable remote-sensing measures can help evaluate public infrastructure policies where administrative road-condition data are unavailable or inconsistent. 

The rest of the paper proceeds as follows. Section \ref{sec:lit_review} positions the study in the literatures on infrastructure, capitalization, and road-quality measurement. Section \ref{sec:data} describes Ohio's local road-finance institutions, the renewal elections, the housing data, and the road-quality measure. Section \ref{sec:empirical} presents the dynamic regression discontinuity design and the identifying assumptions. Section \ref{sec:result} reports house-price, road-quality, heterogeneity, and robustness results. Section \ref{sec:mech} discusses policy mechanisms and interpretation, and Section \ref{sec:conclusion} concludes.
